%% -----------------------------------------------------------
%% Projeto 1 da disciplina de INF2608 - Fundamentos de Computacao Grafica
%% Versao v2 em LaTeX com estrutura orientada a requisitos.

\documentclass[portuguese]{sbcreviews-2025}%
\usepackage[misc,geometry]{ifsym}
\usepackage{aas_macros}
\usepackage[bottom]{footmisc}
\usepackage{tabularray}
\usepackage{afterpage}
\usepackage{url}
\usepackage{pifont}

\setcitestyle{square}
\newcommand{\slidescite}[2]{\citep[pp.~#2]{#1}}
\newcommand{\traceability}[1]{%
    \begin{quote}
    \footnotesize\textbf{Rastreabilidade no codigo.}~#1
    \end{quote}
}

\definecolor{engtitle}{rgb}{0.5,0.5,0.5}
\definecolor{darkblue}{rgb}{0.0,0.0,0.0}
\hypersetup{colorlinks,breaklinks,
            linkcolor=darkblue,urlcolor=darkblue,
            anchorcolor=darkblue,citecolor=darkblue}

\category{Projeto 1 da disciplina de INF2608 -- Fundamentos de Computacao Grafica}
\title[Analise por Requisitos do Projeto 1]{Analise Tecnico-Cientifica por Requisitos de uma Implementacao de Ray Tracing}
\engtitle{\textcolor{engtitle}{Requirement-Oriented Technical Analysis of a Ray Tracing Implementation}}

\author[Yang Ricardo Barcellos Miranda]{
\affil{\textbf{Yang Ricardo Barcellos Miranda}~[\href{mailto:ymiranda@inf.puc-rio.br}{\textit{ymiranda@inf.puc-rio.br}}]}
}

\institution{Pontificia Universidade Catolica do Rio de Janeiro - PUC-Rio}

\begin{document}
\begin{frontmatter}

\maketitle

\begin{abstract-pt}
Este texto apresenta a versao v2 do relatorio do Projeto 1 com organizacao por requisitos do enunciado. O objetivo e tornar a validacao mais direta, mapeando cada requisito principal para uma cena dedicada, um comando de execucao reproduzivel e uma evidencia visual associada, mantendo a mesma configuracao de camera entre os experimentos. A analise preserva a rastreabilidade com o codigo em \path{src/ray_tracing_2/}, destaca aderencia e limites da implementacao atual e separa requisitos principais de extensoes adicionais.
\end{abstract-pt}

\begin{abstract-en}
This v2 report reorganizes the Project 1 analysis around statement requirements. The goal is to provide direct validation by mapping each main requirement to a dedicated scene, a reproducible command-line execution, and associated visual evidence, while keeping camera configuration fixed across experiments. The text preserves code traceability in \path{src/ray_tracing_2/}, highlights adherence and current limitations, and separates core requirements from additional extensions.
\end{abstract-en}

\begin{pchaves}
tracado de raios, computacao grafica, requisitos, Phong, amostragem
\end{pchaves}

\begin{keywords}
ray tracing, computer graphics, requirements, Phong, sampling
\end{keywords}

\begin{dates}
\noindent{\sffamily\textbf{Entrega:}} 10/05/2026
\end{dates}

\end{frontmatter}

\section{Introducao}
\label{sec:introducao-v2}

Esta versao reorganiza o relatorio por requisitos principais de \path{materiais/proj1.pdf}. O principio metodologico e simples: cada requisito deve ser respondido por uma cena principal, um comando reproduzivel e uma interpretacao objetiva do resultado.

\section{Requisitos e Aderencia}
\label{sec:aderencia-v2}

\begin{tblr}{|l|l|l|}
\hline
\textbf{Requisito} & \textbf{Cena principal} & \textbf{Status} \\
\hline
R1: esferas e caixas & \path{proj1_req1_geometry.py} & Implementado \\
\hline
R2: luz pontual & \path{proj1_req2_point_lights.py} & Implementado \\
\hline
R3: Phong + sombras & \path{proj1_req3_phong_shadows.py} & Implementado \\
\hline
R4: multiplas amostras por pixel & \path{proj1_req4_sampling.py} & Implementado \\
\hline
\end{tblr}

\traceability{\path{src/ray_tracing_2/proj1_scene_common.py}, \path{src/ray_tracing_2/proj1_req1_geometry.py}, \path{src/ray_tracing_2/proj1_req2_point_lights.py}, \path{src/ray_tracing_2/proj1_req3_phong_shadows.py}, \path{src/ray_tracing_2/proj1_req4_sampling.py}.}

\section{Procedimento Experimental Padrao}
\label{sec:procedimento-v2}

Os testes de validacao rapida desta versao usam:
\begin{itemize}
    \item resolucao fixa \texttt{800x600};
    \item \texttt{--spp} no intervalo 1 a 4;
    \item baseline em \texttt{--spp 1}.
\end{itemize}

Comandos-base:
\begin{verbatim}
python -m ray_tracing_2.proj1_req1_geometry --width 800 --height 600 --spp 1 --sampling_mode center
python -m ray_tracing_2.proj1_req2_point_lights --width 800 --height 600 --spp 1 --sampling_mode center
python -m ray_tracing_2.proj1_req3_phong_shadows --width 800 --height 600 --spp 1 --sampling_mode center
python -m ray_tracing_2.proj1_req4_sampling --width 800 --height 600 --spp 1 --sampling_mode center
\end{verbatim}

\section{Discussao Tecnica Resumida}
\label{sec:discussao-v2}

\textbf{R1.} A cena principal usa apenas caixas instanciadas e esferas no mesmo enquadramento de camera.

\textbf{R2.} A iluminacao usa exclusivamente fontes pontuais para separar contribuicoes de luz.

\textbf{R3.} O sombreamento local de Phong e as sombras duras aparecem no mesmo setup, com materiais de resposta angular distinta.

\textbf{R4.} O requisito de multiplas amostras e validado por comparacao controlada de \texttt{center} contra modos estocasticos no mesmo quadro.

\section{Limites e Requisitos Adicionais}
\label{sec:limites-v2}

Triangulos, BVH local e luz de area permanecem como extensoes adicionais e nao substituem as evidencias principais dos requisitos 1--4.

\section{Conclusao}
\label{sec:conclusao-v2}

A estrutura v2 melhora a rastreabilidade requisito-a-requisito, mantendo custo de reproducao baixo e coerencia de camera entre cenas.

\bibliographystyle{apalike-sol}
\bibliography{refs}

\end{document}
